% =============================================================================
% 03b-contextos.tex - PARTE 3B, contextos que a matriz nao resolve sozinha.
% Frames 21 e 22. Gabarito fixo: macro unica \quadroDx, geometria identica,
% dois quadros por frame, nenhuma revisao do diagnostico.
%
% Estes dois frames sao obrigatorios. Sem eles o recorte transdiagnostico e
% anunciado e nao entregue, porque os seis casos da atividade nao cobrem
% nenhum destes quatro contextos.
% =============================================================================

% ----------------------------------------------------------------- frame 21 --
\begin{frame}{\ftit{\faUsers}{Contextos que a matriz n\~ao resolve sozinha, I}}
\CRar{0}
\begin{columns}[onlytextwidth, T]
\begin{column}{0.478\textwidth}
\quadroDx{Pessoa idosa}
{Inten\c{c}\~ao mais determinada, menos tentativas e maior letalidade por
tentativa. Doen\c{c}a dolorosa ou incapacitante, perda de autonomia, luto
recente e isolamento.}
{Manejo agressivo da dor e da comorbidade. Envolvimento expl\'icito do
cuidador. Revis\~ao do que est\'a guardado em casa.}
{Normalizar a desesperan\c{c}a como coisa da idade.}
\end{column}
\calha
\begin{column}{0.478\textwidth}
\quadroDx{Adolescente}
{Cont\'agio e efeito de exposi\c{c}\~ao. Humilha\c{c}\~ao, ruptura e
exposi\c{c}\~ao digital. Impulsividade pr\'opria da faixa.}
{Parceria com o respons\'avel, com o adolescente sabendo o que ser\'a dito.
Dever de comunica\c{c}\~ao ao Conselho Tutelar quando couber.}
{Prometer sigilo absoluto antes de saber o que ser\'a contado.}
\end{column}
\end{columns}
\end{frame}

% ----------------------------------------------------------------- frame 22 --
\begin{frame}{\ftit{\faUsers}{Contextos que a matriz n\~ao resolve sozinha, II}}
\CRar{0}
\begin{columns}[onlytextwidth, T]
\begin{column}{0.478\textwidth}
\quadroDx{Transtorno alimentar}
{Na anorexia nervosa, parcela relevante da mortalidade \'e por suic\'idio, e
n\~ao por complica\c{c}\~ao nutricional. Desesperan\c{c}a e vergonha
sustentadas por anos.}
{Perguntar sobre idea\c{c}\~ao em uma consulta que tende a girar s\'o em torno
de peso e de exame laboratorial.}
{Reduzir a avalia\c{c}\~ao ao risco cl\'inico nutricional.}
\end{column}
\calha
\begin{column}{0.478\textwidth}
\quadroDx{Crise psicossocial}
{Perda, viol\^encia, humilha\c{c}\~ao, problema jur\'idico ou financeiro, sem
transtorno estabelecido. Risco real e frequentemente transit\'orio.}
{A conduta \'e companhia e restri\c{c}\~ao de acesso, com reavalia\c{c}\~ao
pr\'oxima enquanto a crise passa.}
{Prescrever no lugar de organizar prote\c{c}\~ao.}
\end{column}
\end{columns}
\end{frame}
